\documentclass[11pt]{article}

\usepackage[margin=1in]{geometry}
\usepackage{amsmath, amssymb}
\usepackage{hyperref}
\usepackage{authblk}

\setlength\parindent{0pt}
\setlength{\parskip}{8pt}

\title{Project Proposal: Microkernel for Machine Learning}

\author{Rohan Yelandur}
\author{Arnav Balaji} 
\author{Matthew Johnson}
\affil{Kernel Sanders (14)}

\date{\today}
\begin{document}

\maketitle

\section{Project Description}

The goal of this project is to design a microkernel for the bfloat16 datatype in BLIS, a datatype commonly used in machine learning settings. We will implement a high performance GEMM algorithm and analyze performance improvements. Our optimized microkernel will be cache aware, use SIMD intrinsics and data packing, and support multiprocessing. The code will be a compatible supplementation to the existing BLIS repository.

\section{Background and Relevance}

Most high performance computing tasks have traditionally relied on 64-bit and 32-bit floating point precision. However, advances in machine learning have created a shift towards data types that take up less memory. Modern machine learning workloads, particularly in deep neural networks, involve large-scale matrix multiplications where reducing memory bandwidth and increasing throughput are critical for performance. Smaller data formats allow for higher throughput with minimal loss of accuracy in the machine learning model. bfloat16 (“Brain Float 16”) was invented by Google’s AI division to be used on their Cloud TPUs \cite{kanwar2019bfloat16}. The problem with standard FP16 is that its number of exponent bits limits the range of values it can represent. bfloat16 uses 8 bits for the exponent and only 7 for the mantissa, maintaining the range of FP32 at a justifiable loss of precision, especially for the problem setting.

High-performance implementations of matrix multiplication are central to both scientific computing and machine learning \cite{Goto2008AnatomyOH}. These implementations are typically based around blocked algorithms that maximize data reuse in cache and minimize memory operations. The core of this is the microkernel: a small, highly optimized routine that performs computation on sub-blocks of matrices using registers and SIMD instructions. The performance of the overall algorithm is largely determined by the efficiency of this microkernel, as it is invoked repeatedly throughout the computation.

The BLIS framework organizes BLAS operations around such microkernels, allowing high performance to be achieved by optimizing a small number of architecture-specific components while keeping the overall structure portable and reusable \cite{Zee2015BLISAF}. Extending this framework to support data types such as bfloat16 is therefore important for enabling high-performance computation in modern machine learning applications.

\section{Expected Results}

\begin{itemize}
    \item \textbf{Implementation:} 
    An optimized microkernel using SIMD intrinsics for bfloat16-based GEMM. 

    \item \textbf{Correctness:} 
    We will validate our implementation by comparing results against a higher-precision reference (FP32 GEMM), as well as a naive 3-loop GEMM implementation. 
    We will report numerical error and ensure correctness for matrices of arbitrary sizes, including those not divisible by block dimensions.

    \item \textbf{Performance:} 
    We will evaluate performance by measuring effective GFLOPS as a function of matrix size. 
\end{itemize}

\section*{Team Plan}

\begin{itemize}
    \item \textbf{Arnav:} Responsible for designing and implementing the microkernel using SIMD intrinsics, including optimizing fused multiply-add (FMA) operations for performance.

    \item \textbf{Matthew:} Responsible for implementing data packing strategies and handling edge cases, ensuring correctness for matrices of arbitrary sizes.

    \item \textbf{Rohan:} Responsible for integrating the implementation into a BLIS-style framework, developing testing infrastructure, and producing benchmarking and analysis documentation.
\end{itemize}


\bibliographystyle{plain}
\bibliography{citations}

\end{document}
